\documentclass[12pt]{article}
\usepackage[margin=1in]{geometry}
\usepackage{enumitem}
\usepackage{titlesec}
\usepackage{parskip}
\usepackage{hyperref}
\usepackage{fontenc}

\title{\textbf{Stokes 2013 Claude Guide}\\
\large Stokes, Susan C., Thad Dunning, Marcelo Nazareno, and Valeria Brusco. \textit{Brokers, Voters, and Clientelism: The Puzzle of Distributive Politics}.\\
New York: Cambridge University Press, 2013.}
\author{}
\date{}

\begin{document}

\maketitle
\tableofcontents
\newpage

%--------------------------------------------------
\section{Central Claims}
%--------------------------------------------------

\begin{itemize}[leftmargin=*, itemsep=4pt]
    \item The book addresses the \textbf{puzzle of distributive politics}: why do parties and machines routinely lavish targeted material benefits on \textbf{loyal, core supporters}---voters whose support is already secure---rather than on \textbf{swing voters}, as canonical formal models (Cox and McCubbins 1986; Dixit and Londregan 1996; Lindbeck and Weibull 1987) predict they should?
    \item The received theoretical anomaly: since a rational vote-maximizing party should spend scarce resources where the electoral return per dollar is highest (i.e., on persuadable swing voters), the empirical regularity that machines target loyalists---documented across Argentina, Venezuela, Mexico, and India---is a genuine puzzle, not just a data problem.
    \item The authors' resolution is a \textbf{theory of broker-mediated distribution}: political machines are not unitary actors but hierarchical organizations composed of \textbf{party leaders} (who control resources and care about aggregate electoral returns) and \textbf{brokers} (local operatives---ward heelers, network leaders, punteros---who possess granular, face-to-face information about individual voters that leaders lack).
    \item Brokers have their own incentives, distinct from and only imperfectly aligned with those of party leaders: brokers want to build and maintain personal followings, extract rents, and minimize the effort and risk involved in monitoring and persuading unfamiliar voters. This creates a classic \textbf{principal--agent (agency) problem}.
    \item Because brokers have better information about \emph{individual} voters than leaders do, but face incentives to economize on effort and to reward personal loyalists, resources are diverted toward \textbf{loyal supporters and members of brokers' own social networks}---not because this is efficient for the party, but because it is easier and more rewarding for the broker. This is \textbf{agency loss} or \textbf{agency slippage}.
    \item Secret ballot means parties and brokers \textbf{cannot directly verify} how any individual voted; clientelism is therefore not simple, contract-enforceable vote buying but a system of \textbf{probabilistic, relational, and repeated exchange} sustained by brokers' local monitoring capacity, social sanctions, and long-term relationships rather than by literal vote verification.
    \item The book identifies (and rejects as insufficient on their own) three alternative explanations for the ``give to loyalists'' anomaly: (1) that loyalty is \textbf{endogenous}---voters become loyal \emph{because} they were rewarded, so apparent core-voter targeting is actually persuasion in disguise; (2) that machines engage in \textbf{turnout buying} rather than vote buying---mobilizing sympathizers who might otherwise abstain, rather than converting opponents; and (3) that \textbf{subcontracting} to brokers, not leaders' own preferences, explains the pattern. The authors show each of these captures part of the story but that broker agency loss is the primary and most general mechanism.
    \item At the aggregate (leader) level, evidence is more mixed and context-dependent: \textbf{party leaders}, unlike brokers, do sometimes target swing districts or respond to programmatic incentives, especially where brokers' organizational reach is weak or where leaders can better monitor broker behavior---producing a \textbf{disjunction between the strategies of leaders and brokers}.
    \item \textbf{Clientelism correlates with poverty}, both across nations (poorer democracies exhibit more nonprogrammatic distribution) and within nations (poorer individuals are more likely to receive targeted handouts). This is explained by the poor's greater sensitivity to small transfers (declining marginal utility of income), by machines' comparative advantage in monitoring and organizing in poor communities, and by weaker programmatic party institutionalization in poorer polities.
    \item At the macro-historical level, the book asks why machine politics and vote buying \textbf{declined and largely disappeared} in today's advanced democracies (Britain, the United States) even though it was pervasive there in the nineteenth century. The authory identify a shift in the \textbf{relative returns} to leaders of programmatic versus clientelist strategies---driven by expanding electorates, secret ballot reform, rising real wages, administrative/bureaucratic capacity, and the changing organizational costs of monitoring---as leaders increasingly find it more profitable to abandon brokers and machines in favor of programmatic appeals.
    \item Normatively, the book asks \textbf{``what's wrong with buying votes?''} Even though clientelism can, in some respects, function as a redistributive or participation-enhancing mechanism, the authors argue it inflicts a \textbf{diversity of harms}---violations of political equality, distortion of voters' autonomous preferences, weakening of programmatic accountability, entrenchment of intermediary power (brokers), and reinforcement of poverty traps---that outweigh the efficiency- or turnout-based arguments sometimes made in its defense.
\end{itemize}

%--------------------------------------------------
\section{Key Concepts and Terms}
%--------------------------------------------------

\begin{itemize}[leftmargin=*, itemsep=4pt]
    \item \textbf{Distributive politics}: The allocation of particularistic material benefits (jobs, cash, goods, favors, local public works) by political actors in exchange for political support, as distinct from programmatic policy that benefits broad categories of citizens according to general rules.
    \item \textbf{Clientelism / nonprogrammatic distribution}: The direct or contingent exchange of material benefits for political support, typically administered face-to-face and targeted at individuals or small groups rather than allocated by impersonal criteria.
    \item \textbf{Programmatic distribution}: Allocation of benefits according to broad, impersonal, rule-based criteria (e.g., universal pensions, means-tested welfare with transparent eligibility rules) not contingent on individual political support.
    \item \textbf{Pork-barrel / particularistic (but not clientelist) distribution}: Geographically targeted benefits (e.g., infrastructure projects for a district) that are not conditioned on the political behavior of individual recipients---distinguished from clientelism proper by the absence of an individual-level contingency.
    \item \textbf{Core voters vs.\ swing voters}: Core (loyal) voters have a high, stable probability of supporting a given party regardless of transfers; swing voters are those whose vote choice is genuinely responsive to marginal incentives. The ``swing voter model'' (Lindbeck-Weibull, Dixit-Londregan) predicts parties target swing voters; the ``core voter model'' (Cox-McCubbins) predicts targeting of loyalists to ensure turnout and reward past support.
    \item \textbf{Brokers}: Local political operatives (also called \emph{punteros} in Argentina) embedded in communities, who have direct, repeated, face-to-face contact with voters and who mediate the distribution of resources from party leaders to individual citizens.
    \item \textbf{Broker-mediated distribution}: The book's central theoretical model---distribution is filtered through brokers whose local informational advantage and personal incentives systematically bias outcomes away from what party leaders alone would choose.
    \item \textbf{Agency loss / agency slippage}: The gap between what a principal (party leader) would want an agent (broker) to do and what the agent actually does, arising from misaligned incentives and the principal's imperfect ability to monitor the agent.
    \item \textbf{Rent-seeking by brokers}: Brokers diverting resources from the party (or from ostensibly-eligible voters) for personal gain, personal networks, or organizational self-maintenance rather than optimal vote-getting.
    \item \textbf{Perceptions of ballot secrecy}: Because parties cannot literally observe individual votes, clientelist exchange depends on voters' \emph{beliefs} about whether their vote can be inferred or monitored (e.g., via turnout records, social observation, or historical practice)---not on actual verification.
    \item \textbf{Turnout buying vs.\ vote buying (persuasion)}: Turnout buying targets sympathizers to ensure they vote; vote buying/persuasion targets opponents or the undecided to convert or demobilize them. The book's Venezuela data are used to disentangle these motives empirically.
    \item \textbf{Machine politics}: An organizational form of party competition built around hierarchies of brokers who distribute particularistic benefits in exchange for votes and mobilization, historically associated with urban political machines (e.g., Tammany Hall) and contemporary developing-democracy party organizations (PJ in Argentina, PRI in Mexico).
    \item \textbf{Iron law of oligarchy / iron law of clientelism (implicit throughout)}: The tendency of any organization reliant on intermediaries to generate entrenched, self-interested intermediary power that resists full control by the nominal leadership---echoing Michels but applied specifically to distributive machines.
\end{itemize}

%--------------------------------------------------
\section{Core Cases and Data}
%--------------------------------------------------

\begin{itemize}[leftmargin=*, itemsep=4pt]
    \item \textbf{Argentina}: The book's primary case, built on original surveys of both \textbf{brokers} (\emph{punteros}) and \textbf{voters} conducted by the authors across multiple provinces (Buenos Aires, C\'ordoba, San Luis, Misiones). Centered on the Peronist party (PJ) machine; used to test the broker-mediated theory directly (Chapters 2 and 4) and to examine subnational government transfers (ATN funds) to municipalities (Chapter 5).
    \item \textbf{Venezuela}: Analyzed using the \textbf{Maisanta database}---a leaked/compiled dataset linking individuals' self-reported political preferences (from the ``\emph{?`Ni\~no o \~Nia?}''/Tasc\'on-Maisanta lists compiled under Ch\'avez) to receipt of \emph{Misiones} social program benefits---allowing the authors to examine targeting by partisan loyalty at the individual level and to distinguish turnout buying from persuasion (Chapter 2, Appendix C).
    \item \textbf{Mexico}: Survey and panel data on the PRI machine's distribution of campaign gifts, used to test whether benefit receipt tracks partisan affect and whether receiving a gift causally shifts opinions of the incumbent party (addressing the endogenous-loyalty explanation) (Chapter 2).
    \item \textbf{India}: Survey data from Bihar, Rajasthan, and Karnataka examining party membership, partisan ties, and benefit receipt, including a survey experiment on partisan ties and benefit receipt (Chapter 2, Appendix D).
    \item \textbf{Cross-national/African and Latin American data}: Used in Chapter 6 to establish the macro-level correlation between national poverty (GDP per capita) and the prevalence of nonprogrammatic distribution, and the individual-level correlation between voter poverty and receipt of targeted rewards.
    \item \textbf{Britain (1832--1923) and the United States (1789--2000)}: The historical core of Chapter 8, using long time series on contested elections, campaign spending, franchise expansion, wages, and urbanization/industrialization to explain the \textbf{decline of vote buying} in both countries---Britain's decline linked to the 1883 Corrupt Practices Act and secret ballot reforms plus rising wages and administrative capacity; the U.S. decline linked to Progressive Era reforms, civil service reform, and the direct primary undermining machine control over nominations.
\end{itemize}

%--------------------------------------------------
\section{Core Works Cited / Engaged}
%--------------------------------------------------

\begin{itemize}[leftmargin=*, itemsep=4pt]
    \item \textbf{Gary Cox and Mathew McCubbins}, ``Electoral Politics as a Redistributive Game'' (1986, \textit{Journal of Politics})---the ``core voter'' model predicting risk-averse parties reward loyal supporters; a central theoretical target that the book's empirics partially vindicate at the broker level.
    \item \textbf{Avinash Dixit and John Londregan}, ``The Determinants of Success of Special Interests in Redistributive Politics'' (1996, \textit{Journal of Politics})---the canonical ``swing voter'' model predicting parties target responsive, contestable voters; the model whose predictions the book's core anomaly most directly challenges.
    \item \textbf{Assar Lindbeck and J\"orgen Weibull}, ``Balanced-Budget Redistribution as the Outcome of Political Competition'' (1987, \textit{Public Choice})---an earlier formalization of swing-voter targeting logic underlying the Dixit-Londregan tradition.
    \item \textbf{Herbert Kitschelt and Steven Wilkinson}, eds., \textit{Patrons, Clients, and Policies} (2007)---the major edited volume establishing the programmatic/clientelistic linkage-strategy distinction that frames the book's conceptual scheme.
    \item \textbf{Susan Stokes}, ``Perverse Accountability: A Formal Model of Machine Politics with Evidence from Argentina'' (2005, \textit{American Political Science Review})---Stokes's own earlier article laying groundwork for the monitoring/perverse-accountability logic developed further here.
    \item \textbf{James Scott}, ``Patron-Client Politics and Political Change in Southeast Asia'' (1972, \textit{APSR})---foundational clientelism scholarship on patron-client dyads as a mode of political organization, an intellectual ancestor of the book's brokerage framework.
    \item \textbf{Robert Michels}, \textit{Political Parties} (1911)---the ``iron law of oligarchy''; relevant to the book's account of how intermediary agents (brokers) accrue independent power within hierarchical organizations.
    \item \textbf{Beatriz Magaloni}, \textit{Voting for Autocracy} (2006)---on the PRI machine in Mexico and the electoral logic of authoritarian clientelism; connects to the book's Mexico data and to the broader literature on machine survival.
    \item \textbf{Steven Wilkinson}, \textit{Votes and Violence} (2004) and work on Indian machine politics---informs the book's India chapters and the comparative brokerage framework.
    \item \textbf{Javier Auyero}, \textit{Poor People's Politics} (2000)---ethnographic work on Peronist clientelist networks in Argentina, providing qualitative texture to the brokerage relationships the book models formally.
    \item \textbf{Frances Rosenbluth and collaborators, and the broader literature on electoral system effects on particularism} (e.g., personal-vote-seeking incentives)---background for why certain electoral rules encourage clientelist versus programmatic strategies.
    \item \textbf{Seymour Martin Lipset (via secondary engagement)} and modernization-theoretic accounts linking development to democratic/programmatic politics---engaged implicitly through the book's poverty--clientelism correlation (Chapter 6) and the historical decline chapters (Chapter 8).
    \item \textbf{Simeon Nichter}, ``Vote Buying or Turnout Buying? Machine Politics and the Secret Ballot'' (2008, \textit{APSR})---directly engaged in Chapter 2's discussion of turnout buying versus vote buying as competing/complementary explanations.
    \item \textbf{Kanchan Chandra}, \textit{Why Ethnic Parties Succeed} (2004)---on patronage and ethnic voting in India, relevant background for the India survey chapters.
\end{itemize}

%--------------------------------------------------
\section{Chapter 1: Between Clients and Citizens---Puzzles and Concepts in the Study of Distributive Politics}
%--------------------------------------------------

\begin{itemize}[leftmargin=*, itemsep=4pt]
    \item \textbf{Framing question}: Why is it legal and normatively acceptable for parties to promise programmatic benefits to broad categories of voters, but illegal and normatively suspect to exchange targeted material benefits for individual votes?
    \item \textbf{Conceptualizing modes of distribution (\S1.1)}: The authors build a typology distinguishing distributive strategies along two axes---(1) whether benefits are targeted at \textbf{individuals/small groups} or at \textbf{broad categories}, and (2) whether receipt is \textbf{contingent} on the recipient's political behavior or not. This produces the core distinction among \textbf{programmatic distribution} (broad, noncontingent), \textbf{pork barrel} (targeted geographically/by group, but noncontingent on individual behavior), and \textbf{clientelism} (targeted at individuals, contingent on their political behavior or perceived loyalty).
    \item \textbf{The conceptual scheme (Figure 1.1)}: Maps existing empirical studies of distributive politics from many countries onto this typology, showing clientelism/nonprogrammatic distribution is empirically widespread, especially (but not exclusively) in developing democracies.
    \item \textbf{Basic questions about distributive politics (\S1.2)}: Who initiates exchange (leaders or brokers)? Who is targeted (loyalists or swing voters, poor or non-poor)? What sustains compliance given the secret ballot? Why does the prevalence of clientelism vary so much across and within countries and over time?
    \item \textbf{Why study clientelism? (\S1.3)}: Clientelism matters because it is empirically pervasive in the developing world, shapes the quality of democratic representation and accountability, has significant welfare/redistributive consequences for the poor, and offers a lens on organizational agency problems within parties more generally. The authors also connect the prevalence of nonprogrammatic distribution to national income (Figure 1.2 shows GDP per capita trends distinguishing democracies and autocracies), foreshadowing the poverty argument developed later.
    \item \textbf{Structure of the book (\S1.4)}: Lays out the book's three-part architecture---Part II (``the micro-logic of clientelism,'' Chapters 2--6) develops and tests the broker-mediated theory; Part III (``the macro-logic of vote buying,'' Chapters 7--8) explains cross-national and historical variation in the prevalence of machine politics; Part IV (Chapter 9) turns to normative democratic theory.
    \item \textbf{A comment on research methods (\S1.5)}: The book self-consciously combines original survey research (voters and brokers), secondary/comparative statistical analysis, and historical narrative---an explicitly multi-method, triangulating design chosen because no single data source can resolve the ``who is targeted and why'' puzzle on its own.
\end{itemize}

%--------------------------------------------------
\section{Chapter 2: Gaps Between Theory and Fact}
%--------------------------------------------------

\begin{itemize}[leftmargin=*, itemsep=4pt]
    \item \textbf{Purpose}: Lays out the central empirical anomaly in detail---across multiple countries and datasets, benefits go disproportionately to \textbf{loyal partisans}, not to swing/persuadable voters, contradicting the dominant swing-voter formal models (\S2.1, \textit{Theories of Distributive Politics}).
    \item \textbf{Evidence assembled}: Original and secondary survey data from Argentina, Venezuela, Mexico, and India (Table 2.1 catalogs primary survey data and sources) consistently show reward recipients are disproportionately strong partisans of the distributing party, not fence-sitters (e.g., Argentina Figure 2.2--2.3 on Peronist loyalists; Venezuela Figures 2.4--2.7 using the Maisanta database on \emph{Misiones} beneficiaries; Mexico Figure 2.8 on PRI gift receipt by partisanship; India Table 2.3 on party membership and benefit receipt).
    \item \textbf{Explaining the anomaly---is ``loyalty'' endogenous? (\S2.2)}: Considers the possibility that apparent core-voter targeting is illusory---voters \emph{become} loyal \emph{because} of past receipt of benefits (reverse causality/endogeneity). Using panel data (e.g., Mexico 2000 panel, Figure 2.10--2.11) and instrumental-variable regressions (Argentina, Table 2.4), the authors find some support for an endogenous-loyalty/persuasion effect but conclude it does not fully explain the targeting pattern---machines still disproportionately \emph{target} pre-existing loyalists rather than genuine swing voters.
    \item \textbf{Explaining the anomaly---turnout buying? (\S2.3)}: Tests whether rewards go to loyalists because parties are trying to \textbf{mobilize turnout among sympathizers} rather than to persuade opponents (Table 2.6 on electoral rewards and two-dimensional voter types; Table 2.7 and Figure 2.12 on persuasion versus mobilization in Venezuela). Finds real evidence of turnout buying, especially among loyal-but-otherwise-unlikely-to-vote citizens, but again only a partial explanation.
    \item \textbf{Explaining the anomaly---subcontracting? (\S2.4)}: Considers whether leaders deliberately \textbf{delegate} distribution to brokers knowing this will produce loyalist-biased outcomes (i.e., the pattern reflects leader strategy, not broker agency loss). The chapter sets up---but does not fully resolve---this question, motivating the formal theory of broker-mediated distribution developed in Chapter 3.
    \item \textbf{Conclusion of the chapter}: None of the three ``patches'' to the swing-voter model (endogenous loyalty, turnout buying, subcontracting) alone accounts for the full empirical pattern; a more fundamental reworking of the theory---centered on brokers as independent strategic actors---is required.
\end{itemize}

%--------------------------------------------------
\section{Chapter 3: A Theory of Broker-Mediated Distribution}
%--------------------------------------------------

\begin{itemize}[leftmargin=*, itemsep=4pt]
    \item \textbf{Core theoretical move}: Decomposes the ``party'' into two strategic levels---\textbf{party leaders}, who control aggregate resources and care about vote share/seats, and \textbf{brokers}, who have local, granular information about individual voters but their own distinct payoffs.
    \item \textbf{A model of rent-seeking brokers (\S3.1)}: Formalizes brokers as agents who can \textbf{divert} (``rent-seek'' from) party resources intended for voters, or who allocate benefits to build/maintain their own personal following rather than to maximize the party's vote return. Because leaders cannot perfectly observe how brokers behave (imperfect monitoring), brokers have latitude to shade allocations toward their own interests---typically toward voters who are already loyal to the broker personally, who are cheap to monitor, or who reinforce the broker's own local standing.
    \item \textbf{The objectives of party leaders (\S3.2)}: Leaders are modeled as caring about the party's aggregate electoral return and are willing to trade off some efficiency for organizational control, but face a fundamental information asymmetry: only brokers know individual voters' preferences and responsiveness with precision.
    \item \textbf{The implications of agency loss (\S3.3)}: Because brokers are rewarded for maintaining networks (not, strictly, for maximizing marginal votes for the party), the model predicts systematic targeting toward loyalists and personal clients of brokers---generating, as an equilibrium outcome, precisely the loyalist-biased distribution pattern documented empirically in Chapter 2, \emph{without} needing to assume irrationality on the part of leaders, voters, or brokers.
    \item \textbf{Significance}: This is the book's central theoretical contribution---it reframes the ``distributive politics puzzle'' not as a question about voter psychology or about which single unified actor targets whom, but as a \textbf{principal--agent problem internal to political organizations}, shifting the unit of analysis from the party as a unitary actor to the party as a hierarchy of differently-incentivized actors.
\end{itemize}

%--------------------------------------------------
\section{Chapter 4: Testing the Theory of Broker-Mediated Distribution}
%--------------------------------------------------

\begin{itemize}[leftmargin=*, itemsep=4pt]
    \item \textbf{Data source}: The authors' original \textbf{Argentina Brokers' Survey} (detailed in Appendix A), covering hundreds of \emph{punteros} across multiple provinces, is the empirical backbone of this chapter.
    \item \textbf{Who are the brokers? (\S4.1)}: Characterizes brokers demographically and organizationally---typically locally embedded party activists (elected councilors and non-elected operatives) who maintain personal networks of clients (Table A.2 breaks down sampled brokers by province).
    \item \textbf{Can brokers infer voters' choices? (Figure 4.1)}: Tests whether brokers claim (and plausibly have) the capacity to infer how specific voters will vote or did vote---central to the theory's assumption of an informational advantage over leaders, despite the secret ballot.
    \item \textbf{Monitoring capacity and ``hard to get'' voters (Figure 4.2--4.3)}: Explores how many voters resist clear signaling of their loyalty (``play hard to get'') and whether brokers can identify them, plus the extent of benefit offers versus voter-initiated requests for help.
    \item \textbf{Brokers' preference for loyal voters (Figure 4.4, survey experiment)}: A direct experimental test showing brokers, when given discretion, systematically prefer to reward already-loyal clients over uncertain or opposition-leaning voters---directly consistent with the agency-loss prediction.
    \item \textbf{Network structure (Figures 4.5--4.7)}: Documents the size and composition of brokers' personal networks, including the proportion of clients who were already party sympathizers before being recruited into the broker's network, and the relationship between rally attendance/network participation and reported voting behavior for the broker's candidate.
    \item \textbf{Rent-seeking and organization-building by brokers (Figures 4.8--4.9)}: Direct survey evidence that a meaningful share of brokers keep some party-allocated benefits for themselves (rent-seeking) or condition distribution on rally attendance and organizational participation rather than strictly on electoral responsiveness (organization-building)---both patterns consistent with agency loss rather than pure vote-maximization.
    \item \textbf{Monitoring and sanctioning of brokers by leaders (Figures 4.10--4.11, Table 4.1)}: Examines how party leaders evaluate broker performance (what criteria they use) and how brokers perceive their own vulnerability to having resources or their position withdrawn---evidence of leaders' (imperfect) attempts to counteract agency loss through incentive and monitoring mechanisms.
    \item \textbf{Conclusion (\S4.3)}: The broker survey evidence strongly supports the theory's core assumptions and predictions: brokers do have privileged local information, do exercise real discretion, and do systematically favor loyalists and personal network members, generating agency loss that is only partially checked by leaders' monitoring efforts.
\end{itemize}

%--------------------------------------------------
\section{Chapter 5: A Disjunction Between the Strategies of Leaders and Brokers?}
%--------------------------------------------------

\begin{itemize}[leftmargin=*, itemsep=4pt]
    \item \textbf{Purpose}: Shifts the level of analysis upward, asking whether \textbf{party leaders}---as opposed to brokers---behave differently, targeting swing districts or responding more to programmatic/aggregate electoral logic when they retain direct control over distribution (\S5.1, \textit{Theories of Distribution by Party Leaders}).
    \item \textbf{Do swing districts receive party largess? (\S5.2)}: Reviews and reanalyzes ecological (district-level) studies of distributive politics (Table 5.2 codes many such studies by whether leaders' strategy resembles core-voter or swing-voter targeting), alongside original terminology distinguishing types of subnational districts by their partisan competitiveness (Table 5.1).
    \item \textbf{Leaders and brokers in four developing democracies (\S5.3)}: Compares evidence across Argentina, Venezuela, Mexico, and India to assess whether the leader-level pattern diverges from the broker-level pattern documented in Chapters 2 and 4. Uses the Argentine case of \textbf{ATN (Aportes del Tesoro Nacional) discretionary transfers} to municipalities (Figure 5.1, Table 5.3) as a concrete test of whether higher-level distribution (controlled more directly by leaders/governors) follows swing-district logic more than broker-level distribution does.
    \item \textbf{Key finding}: There is indeed a partial \textbf{disjunction}---leaders, especially when they can bypass or closely monitor brokers, show more evidence of targeting political competitiveness (swing logic) than brokers do at the individual level, where loyalist-targeting dominates. This supports the book's core claim that the ``who gets targeted'' puzzle is partly an artifact of \emph{which level of the organizational hierarchy} is doing the targeting.
    \item \textbf{Conclusion (\S5.4)}: The apparent contradiction between theory (predicting swing-voter targeting) and fact (loyalist targeting) is substantially resolved once distributive politics is disaggregated by organizational level---reinforcing the broker-mediated framework as more than just an Argentina-specific curiosity.
\end{itemize}

%--------------------------------------------------
\section{Chapter 6: Clientelism and Poverty}
%--------------------------------------------------

\begin{itemize}[leftmargin=*, itemsep=4pt]
    \item \textbf{Framing (\S6.1)}: Distinguishes two related but distinct empirical regularities---the \textbf{poverty of nations} (poorer countries exhibit more nonprogrammatic distribution) and the \textbf{poverty of voters} (within a country, poorer individuals are more likely to receive clientelist rewards).
    \item \textbf{National poverty and nonprogrammatic distribution (\S6.2)}: Cross-national evidence (Figures 6.1--6.3, covering Africa and Latin America) shows a clear negative relationship between GDP per capita and the prevalence of campaign gift-giving/clientelist distribution, both across countries and, in the Latin American 2010 data, in a more fine-grained way.
    \item \textbf{Individual poverty and nonprogrammatic distribution (\S6.3)}: Within-country survey evidence from Argentina and Venezuela (Figures 6.4--6.5) shows poorer individuals disproportionately report receiving targeted rewards.
    \item \textbf{Why do machines target the poor? (\S6.4)}: Offers several complementary mechanisms: (1) \textbf{declining marginal utility of income}---small transfers matter more to poor voters, making clientelist rewards more effective persuasion/mobilization tools among the poor; (2) machines' \textbf{comparative organizational advantage} in monitoring and mobilizing in poor, often geographically concentrated communities; (3) the poor's greater reliance on informal/patronage-based welfare provision where state programmatic capacity is weak; and evidence (Figure 6.6, African survey data) that poor respondents are more likely to believe politicians' promises, making them more responsive to clientelist appeals generally.
    \item \textbf{Significance}: Connects the book's micro-level (broker) theory to a macro-level, cross-nationally robust empirical regularity, and lays groundwork for Chapters 7--8's account of why rising development/wages undermines machine politics.
\end{itemize}

%--------------------------------------------------
\section{Chapter 7: Party Leaders Against the Machine}
%--------------------------------------------------

\begin{itemize}[leftmargin=*, itemsep=4pt]
    \item \textbf{Transition to Part III}: Having established the micro-logic of broker-mediated clientelism, the book turns to the ``macro-logic'' question: what explains variation over time and across countries in whether party leaders choose to \textbf{rely on machines} at all, versus shifting toward \textbf{programmatic} politics?
    \item \textbf{Broker-mediated theory and the returns to clientelism (\S7.1)}: Argues that because brokers extract agency rents, clientelism is costly and inefficient \emph{for leaders}---leaders would prefer programmatic strategies when the relative returns favor them, i.e., machine politics persists only where its benefits to leaders (mobilization, information, organizational control) still outweigh the agency losses documented in Chapters 3--4.
    \item \textbf{Clientelism and programmatic politics: a model (\S7.2)}: Formalizes leaders' choice between investing in machine infrastructure (brokers) versus programmatic appeals as a function of the relative costs and returns of each strategy, which in turn depend on features like electorate size, monitoring costs, and the developmental/administrative context.
    \item \textbf{When do leaders choose machine politics? (\S7.3)}: Identifies conditions favoring machine politics: relatively small, geographically concentrated, or poor electorates (lower monitoring/organizing costs and higher marginal value of transfers); weak state administrative/bureaucratic capacity (making programmatic delivery costly or noncredible); and weak alternative sources of information about voters' preferences.
    \item \textbf{Testing the theory (\S7.4)}: Draws on cross-national and historical evidence (anticipating Chapter 8's British and U.S.\ material as well as contemporary comparisons) to show that shifts in these conditions---rising incomes, growing and dispersed electorates, state-building, secret ballot and civil-service reforms---predict the shift away from machine politics (Figure 7.1 summarizes the factors encouraging the shift to programmatic politics).
    \item \textbf{Significance}: Provides the theoretical bridge between the book's micro-level broker theory and its historical account (Chapter 8) of vote buying's decline in advanced democracies---clientelism is not a fixed cultural trait but an organizational strategy whose relative profitability to leaders changes with development.
\end{itemize}

%--------------------------------------------------
\section{Chapter 8: What Killed Vote Buying in Britain and the United States?}
%--------------------------------------------------

\begin{itemize}[leftmargin=*, itemsep=4pt]
    \item \textbf{Purpose (\S8.1)}: Applies the book's theory of leaders' strategic choice (Chapter 7) to explain the historical \textbf{decline of vote buying and machine politics} in two paradigmatic advanced democracies where clientelism was once pervasive.
    \item \textbf{Britain (\S8.2)}: Uses long time series (1832--1923) on election petitions challenging corrupt/bribed elections (Figures 8.1--8.2), campaign spending on agents and printing (Figure 8.5), population and franchise growth (Figures 8.6--8.8), campaign expenditures per voter (Figure 8.9), and socioeconomic indicators---real wages (Figure 8.10) and the shift from agriculture to industry (Figure 8.11)---to show that vote buying declined as the electorate expanded (making individual-level bribery prohibitively costly per vote), as the 1883 Corrupt and Illegal Practices Prevention Act and earlier secret ballot reform (1872) raised the legal and informational costs of vote buying, and as rising wages reduced the poor's susceptibility to small material inducements.
    \item \textbf{The United States (\S8.3)}: Uses parallel time series (1789--2000) on contested congressional elections (Figures 8.3--8.4) to trace the decline of American urban machine politics (e.g., Tammany Hall), attributing it to the \textbf{Australian (secret) ballot reform} adopted state-by-state beginning in the 1880s--90s, Progressive Era \textbf{civil service reform} (undermining patronage-based machine employment), and the spread of the \textbf{direct primary} (which reduced machines' control over nominations and thus over the loyalty incentives brokers could offer).
    \item \textbf{Comparative synthesis (Figures 8.12--8.14)}: Directly compares population growth, votes cast, and turnout trends in Britain and the United States over the same period, showing broadly parallel declines in vote buying despite different institutional paths, reinforcing the general theory (Chapter 7) over country-specific idiosyncratic explanations.
    \item \textbf{Conclusion (\S8.4)}: The decline of machine politics in both cases is best explained not by a single reform but by the \textbf{joint effect of secret ballot, expanding and increasingly costly-to-monitor electorates, administrative/bureaucratic reform, and rising real incomes}---all of which raised the relative cost, and lowered the relative return, of clientelist over programmatic strategies for party leaders, exactly as the Chapter 7 model predicts.
\end{itemize}

%--------------------------------------------------
\section{Chapter 9: What's Wrong with Buying Votes?}
%--------------------------------------------------

\begin{itemize}[leftmargin=*, itemsep=4pt]
    \item \textbf{Transition to Part IV}: Having developed and tested a positive (explanatory) theory of clientelism, the book closes with a normative assessment grounded in democratic theory.
    \item \textbf{Distributive politics and democratic theory (\S9.1)}: Frames the normative stakes---nonprogrammatic distribution implicates core democratic values including political equality, voter autonomy, and the quality of electoral accountability.
    \item \textbf{Nonprogrammatic distribution and the diversity of harms (\S9.2, Table 9.1 summarizes harms)}: Identifies multiple distinct harms rather than a single unified objection: (1) violation of \textbf{political equality}, since clientelism channels political influence through material need and network access rather than equal citizenship; (2) \textbf{distortion of preference expression}, since votes are induced by side payments rather than reflecting policy views; (3) weakened \textbf{programmatic accountability}, since politicians who can win via clientelism face less pressure to deliver broad public goods or sound policy; (4) \textbf{entrenchment of intermediary (broker) power}, creating durable local political bosses with interests independent of both voters and party leaders; and (5) potential \textbf{reinforcement of poverty}, if clientelism substitutes for programmatic antipoverty policy or if it is used to discourage recipients' independent political and economic development.
    \item \textbf{Arguments in favor of nonprogrammatic distribution (\S9.3)}: Engages seriously with counterarguments---that clientelism can function as an informal, flexible \textbf{redistributive mechanism} reaching the poor where formal welfare states are weak or absent; that it can \textbf{increase political participation and turnout} among otherwise marginalized citizens; and that it can provide efficient, locally-tailored responses to need via brokers' superior local information. The authors take these arguments seriously but ultimately find them insufficient to offset the harms identified in \S9.2, especially given the demonstrated agency losses (Chapters 3--4) that mean clientelist resources are frequently \emph{not} allocated efficiently even by the redistributive or participatory standards its defenders invoke.
    \item \textbf{Conclusions (\S9.4)}: Clientelism is best understood not as a benign alternative development path but as a \textbf{transitional and normatively costly} mode of distributive politics that societies tend to abandon as administrative capacity, income, and electoral institutions develop (per Chapters 6--8)---and that, while not without some redeeming functions, imposes real and multifaceted harms on democratic quality that justify continued reform efforts (ballot secrecy enforcement, bureaucratic professionalization, programmatic party-building) aimed at reducing its prevalence.
\end{itemize}

%--------------------------------------------------
\section{Appendices (Methodological Notes)}
%--------------------------------------------------

\begin{itemize}[leftmargin=*, itemsep=4pt]
    \item \textbf{Appendix A---Argentina Brokers' Survey}: Details the sampling design (Figure A.1), sampled municipalities by province (Table A.1), the number of councilor and non-elected brokers sampled (Table A.2), and survey completion rates by province/municipality (Tables A.3.1--A.3.2)---the empirical foundation for Chapter 4.
    \item \textbf{Appendix B---Argentina Voters' Surveys}: Documents the voter-side survey instruments used alongside the brokers' survey, underlying the Argentina evidence in Chapters 2, 5, and 6.
    \item \textbf{Appendix C---Venezuela Voters' Survey and the Maisanta Database}: Describes the \textbf{Maisanta} database (built from the Ch\'avez government's Tasc\'on list linking citizens' 2004 recall-referendum signatures/preferences to later benefit receipt), including its variables (Table C.1), covariate balance tests (Table C.2), and analysis of whether ideology predicts missing data (Table C.3)---plus contextual data on Ch\'avez approval ratings (Figure C.1). This is among the book's most methodologically novel and ethically sensitive data sources, given its origins in what critics describe as a politically weaponized surveillance list.
    \item \textbf{Appendix D---India Voters' Survey}: Documents the representativeness of the regression-discontinuity study group used in the Karnataka survey (Table D.1), underlying the India evidence in Chapters 2 and 4.
\end{itemize}

\end{document}
